\documentclass[11pt]{article}
\usepackage{amsmath,amssymb,amsthm}
\usepackage{algorithm}
\usepackage{algorithmic}
\usepackage{enumitem}
\usepackage{hyperref}
\usepackage{geometry}
\geometry{margin=1in}
\newtheorem{theorem}{Theorem}
\newtheorem{lemma}{Lemma}
\newtheorem{assumption}{Assumption}
\newcommand{\E}{\mathbb{E}}
\newcommand{\R}{\mathbb{R}}

\begin{document}

%%%%%%%%%%%%%%%%%%%%%%%%%%%%%%%%%%%%%%%%%%%%%%%%%%%%%%%%%%%%%%%%%%%%%%%%%%%%
\section*{Clipped Stochastic Gradient Descent (Clipped SGD)}
%%%%%%%%%%%%%%%%%%%%%%%%%%%%%%%%%%%%%%%%%%%%%%%%%%%%%%%%%%%%%%%%%%%%%%%%%%%%

\section*{Mathematical Formulation}
Let $f:\R^{d}\rightarrow\R$ be differentiable.  
At iteration $t$ we have a (possibly stochastic) gradient estimator 
\[
\tilde g_t:=\nabla f(x_t;\xi_t),
\]
where $\xi_t$ denotes the random data sample (or noise) at time $t$.  
Define the \emph{clipping operator} with threshold $C>0$ as
\[
\operatorname{Clip}_C(v):=
\frac{v}{\max\!\bigl(1,\; \|v\|/C\bigr)}=
\begin{cases}
v, &\|v\|\le C,\\[4pt]
C\,\dfrac{v}{\|v\|}, &\|v\|>C .
\end{cases}
\tag{1}
\]
The clipped gradient is
\[
g_t:=\operatorname{Clip}_C(\tilde g_t). \tag{2}
\]
The update rule of Clipped SGD is
\[
x_{t+1}=x_t-\eta_t\, g_t, \qquad t=0,1,\dots ,T-1, \tag{3}
\]
with a stepsize (learning rate) sequence $\{\eta_t\}_{t\ge0}$.

\section*{Pseudocode}
\begin{algorithm}[H]
\caption{Clipped Stochastic Gradient Descent}
\begin{algorithmic}[1]
\REQUIRE Initial point $x_0\in\R^d$, stepsize schedule $\{\eta_t\}$, clipping threshold $C>0$.
\FOR{$t=0$ \TO $T-1$}
    \STATE Sample a random minibatch (or data point) $\xi_t$.
    \STATE Compute raw stochastic gradient $\tilde g_t=\nabla f(x_t;\xi_t)$.
    \STATE \textbf{Clip:} $g_t \gets \operatorname{Clip}_C(\tilde g_t)$ using (1).
    \STATE Update: $x_{t+1}\gets x_t-\eta_t g_t$.
\ENDFOR
\RETURN $x_T$ (or the averaged iterate $\bar x_T:=\frac1T\sum_{t=0}^{T-1}x_t$).
\end{algorithmic}
\end{algorithm}

\section*{Dry Run Example}
Consider the one‑dimensional quadratic $f(w)=(w-3)^2$.  
Its gradient is $\nabla f(w)=2(w-3)$.  
Take $w_0=0$, stepsize $\eta=0.1$, clipping threshold $C=2$ and **deterministic** gradients (no stochasticity).

\begin{center}
\begin{tabular}{c|c|c|c|c}
$t$ & $w_t$ & $\nabla f(w_t)$ & $g_t=\operatorname{Clip}_C(\nabla f(w_t))$ & $f(w_t)$ \\ \hline
0 & $0$ & $-6$ & $\displaystyle 2\frac{-6}{6}= -2$ & $9$\\
1 & $0-0.1(-2)=0.2$ & $-5.6$ & $-2$ (clipped) & $(0.2-3)^2=7.84$\\
2 & $0.2-0.1(-2)=0.4$ & $-5.2$ & $-2$ (clipped) & $(0.4-3)^2=6.76$\\
3 & $0.4-0.1(-2)=0.6$ & $-4.8$ & $-2$ (clipped) & $(0.6-3)^2=5.76$
\end{tabular}
\end{center}
The clipping keeps the update magnitude at $2$, preventing the huge raw gradient $-6$ from moving the iterate too far.

%%%%%%%%%%%%%%%%%%%%%%%%%%%%%%%%%%%%%%%%%%%%%%%%%%%%%%%%%%%%%%%%%%%%%%%%%%%%
\section*{Assumptions}
%%%%%%%%%%%%%%%%%%%%%%%%%%%%%%%%%%%%%%%%%%%%%%%%%%%%%%%%%%%%%%%%%%%%%%%%%%%%
\begin{assumption}[Smoothness]\label{ass:smooth}
$f$ is $L$‑smooth, i.e.
\[
\|\nabla f(x)-\nabla f(y)\|\le L\|x-y\|, \quad\forall x,y\in\R^{d}.
\tag{A1}
\]
Equivalently,
\[
f(y)\le f(x)+\langle\nabla f(x),y-x\rangle+\frac{L}{2}\|y-x\|^{2}.
\tag{A1$'$}
\]
\end{assumption}

\begin{assumption}[Lower Boundedness]\label{ass:lb}
$f$ is bounded below: $f^*:=\inf_{x}f(x)>-\infty$.
\tag{A2}
\end{assumption}

\begin{assumption}[Unbiased Stochastic Gradient]\label{ass:unbiased}
For every $t$,
\[
\E_{\xi_t}\big[\tilde g_t\mid x_t\big]=\nabla f(x_t),\qquad
\E_{\xi_t}\big[\|\tilde g_t-\nabla f(x_t)\|^{2}\mid x_t\big]\le\sigma^{2}.
\tag{A3}
\]
\end{assumption}

\begin{assumption}[Stepsize]\label{ass:stepsize}
The deterministic stepsize satisfies
\[
0<\eta_t\le\frac{1}{L},\qquad \eta_t\equiv\eta\ \text{(constant)}\ \text{or}\ \eta_t=\frac{c}{\sqrt{t+1}}.
\tag{A4}
\]
\end{assumption}

\begin{assumption}[Clipping Threshold]\label{ass:clip}
The clipping constant $C>0$ is finite and independent of $t$.
\tag{A5}
\end{assumption}

%%%%%%%%%%%%%%%%%%%%%%%%%%%%%%%%%%%%%%%%%%%%%%%%%%%%%%%%%%%%%%%%%%%%%%%%%%%%
\section*{Main Convergence Theorem}
%%%%%%%%%%%%%%%%%%%%%%%%%%%%%%%%%%%%%%%%%%%%%%%%%%%%%%%%%%%%%%%%%%%%%%%%%%%%
\begin{theorem}[Convergence of Clipped SGD]\label{thm:main}
Assume \ref{ass:smooth}–\ref{ass:clip} hold and let $x_{t+1}=x_t-\eta g_t$ with $g_t$ defined in (2).  
Define the averaged squared gradient norm
\[
\Phi_T:=\frac{1}{T}\sum_{t=0}^{T-1}\E\big[\|\nabla f(x_t)\|^{2}\big].
\]
If the constant stepsize satisfies $0<\eta\le\frac{1}{L}$, then for any $T\ge1$
\[
\Phi_T\ \le\ \frac{2\big(f(x_0)-f^*\big)}{\eta T}
\;+\;2L\,C^{2}\eta\;+\;\frac{2\sigma^{2}}{C^{2}}\,\eta .
\tag{4}
\]
Consequently, choosing $\eta=\Theta\!\big(\frac{1}{\sqrt{T}}\big)$ yields
\[
\Phi_T = O\!\Big(\frac{1}{\sqrt{T}}\Big).
\]
\end{theorem}

\section*{Theoretical Properties}
\begin{enumerate}[label=\arabic*.]
\item \textbf{Stability.} The clipping operator guarantees $\|g_t\|\le C$ for all $t$, which together with $\eta\le 1/L$ prevents the iterates from exploding even when raw stochastic gradients have unbounded variance.
\item \textbf{Bias–Variance Trade‑off.} Clipping introduces a deterministic bias
\[
b_t:=\E[g_t\mid x_t]-\nabla f(x_t),
\]
satisfying $\|b_t\|\le \frac{\E\big[(\|\tilde g_t\|-C)_{+}^{2}\big]}{C}\le \frac{\sigma^{2}}{C^{2}}$ (Lemma~\ref{lem:bias}). The bound appears as the third term in (4).
\item \textbf{Comparison to vanilla SGD.} When $C\ge \max_t\|\tilde g_t\|$ the clipping is never active, $g_t=\tilde g_t$, and (4) reduces to the standard SGD rate $O(1/\sqrt{T})$ (up to the $2LC^{2}\eta$ term which vanishes as $C\to\infty$).
\item \textbf{Hyper‑parameter sensitivity.} A larger $C$ reduces bias but increases the $2LC^{2}\eta$ term; a smaller $C$ improves stability at the cost of larger bias. The optimal $C$ balances both contributions, typically $C=\Theta(\sigma^{1/2}L^{-1/4})$ for the best $O(T^{-1/2})$ constant.
\end{enumerate}

%%%%%%%%%%%%%%%%%%%%%%%%%%%%%%%%%%%%%%%%%%%%%%%%%%%%%%%%%%%%%%%%%%%%%%%%%%%%
\section*{Preliminaries (Definitions and Lemmas)}
%%%%%%%%%%%%%%%%%%%%%%%%%%%%%%%%%%%%%%%%%%%%%%%%%%%%%%%%%%%%%%%%%%%%%%%%%%%%

\begin{lemma}[Descent Lemma]\label{lem:descent}
If $f$ is $L$‑smooth then for any $x$ and any vector $d$,
\[
f(x-d)\le f(x)-\langle\nabla f(x),d\rangle+\frac{L}{2}\|d\|^{2}.
\tag{5}
\]
\end{lemma}
\begin{proof}
Directly from \ref{ass:smooth} (A1$'$) with $y=x-d$.
\end{proof}

\begin{lemma}[Properties of the clipping operator]\label{lem:clip}
For any $v\in\R^{d}$ and $C>0$ let $c(v):=\operatorname{Clip}_C(v)$.
\begin{enumerate}[label=(\alph*)]
\item $\|c(v)\|\le C$ and $c(v)$ is colinear with $v$.
\item $\langle c(v),v\rangle = \|c(v)\|\,\|v\| = \min\{C,\|v\|\}\,\|v\|$.
\item $\|c(v)-v\|\le\big(\|v\|-C\big)_{+}\le\frac{\|v\|^{2}}{C}$, where $(a)_{+}:=\max\{a,0\}$.
\end{enumerate}
\end{lemma}
\begin{proof}
From definition (1). If $\|v\|\le C$, $c(v)=v$ and all statements are trivial.  
If $\|v\|>C$, then $c(v)=C\,v/\|v\|$, whence $\|c(v)\|=C$, colinearity, and
\[
\langle c(v),v\rangle = C\frac{\langle v,v\rangle}{\|v\|}=C\|v\|.
\]
Moreover,
\[
\|c(v)-v\| = \Big\|C\frac{v}{\|v\|}-v\Big\|
= \Big|\,\|v\|-C\,\Big| = \|v\|-C,
\]
which proves (c). The last inequality follows from $\|v\|-C\le \|v\|^{2}/C$ for $\|v\|\ge C$.
\end{proof}

\begin{lemma}[Bias bound]\label{lem:bias}
Let $g_t=\operatorname{Clip}_C(\tilde g_t)$ with $\tilde g_t$ satisfying (A3).  
Define the conditional bias $b_t:=\E[g_t\mid x_t]-\nabla f(x_t)$. Then
\[
\|b_t\|\le \frac{\sigma^{2}}{C^{2}}.
\tag{6}
\]
\end{lemma}
\begin{proof}
From Lemma~\ref{lem:clip} (c),
\[
\|g_t-\tilde g_t\|\le\frac{\|\tilde g_t\|^{2}}{C}.
\]
Take conditional expectation and use $\E[\|\tilde g_t\|^{2}\mid x_t]
= \|\nabla f(x_t)\|^{2}+ \E[\|\tilde g_t-\nabla f(x_t)\|^{2}\mid x_t]
\le \|\nabla f(x_t)\|^{2}+\sigma^{2}$.
Hence
\[
\|b_t\|
= \big\|\E[g_t-\tilde g_t\mid x_t]\big\|
\le \frac{ \|\nabla f(x_t)\|^{2}+\sigma^{2}}{C}
\le \frac{\sigma^{2}}{C^{2}}+\frac{\|\nabla f(x_t)\|^{2}}{C}.
\]
Dropping the non‑negative term $\|\nabla f(x_t)\|^{2}/C$ yields (6).
\end{proof}

%%%%%%%%%%%%%%%%%%%%%%%%%%%%%%%%%%%%%%%%%%%%%%%%%%%%%%%%%%%%%%%%%%%%%%%%%%%%
\section*{Main Proof (Step‑by‑Step Derivation)}
%%%%%%%%%%%%%%%%%%%%%%%%%%%%%%%%%%%%%%%%%%%%%%%%%%%%%%%%%%%%%%%%%%%%%%%%%%%%
We prove Theorem~\ref{thm:main}.  
Throughout the proof we keep the notation $x_t$ for the iterate, $g_t$ for the clipped stochastic gradient, and $d_t:=\eta g_t$ for the actual step.

\subsection*{Step 1 – Apply the descent lemma}
From Lemma~\ref{lem:descent} with $d_t=\eta g_t$,
\[
f(x_{t+1})\le f(x_t)-\eta\langle\nabla f(x_t),g_t\rangle
+\frac{L\eta^{2}}{2}\|g_t\|^{2}.
\tag{7}
\]

\subsection*{Step 2 – Replace the inner product using Lemma~\ref{lem:clip}}
By Lemma~\ref{lem:clip} (b),
\[
\langle\nabla f(x_t),g_t\rangle
= \min\big\{C,\|\nabla f(x_t)\|\big\}\,\|\nabla f(x_t)\|.
\tag{8}
\]
Also $\|g_t\|=\min\{C,\|\tilde g_t\|\}\le C$.

\subsection*{Step 3 – Decompose the stochasticity}
Write $g_t = \tilde g_t - (\tilde g_t-g_t)$.  
Take conditional expectation w.r.t.\ $\xi_t$ (equivalently w.r.t.\ $g_t$) while keeping $x_t$ fixed.
Using unbiasedness (A3) and the bias bound (Lemma~\ref{lem:bias}),
\[
\E\big[\langle\nabla f(x_t),g_t\rangle\mid x_t\big]
= \|\nabla f(x_t)\|^{2}+ \langle\nabla f(x_t),b_t\rangle
\ge \|\nabla f(x_t)\|^{2} - \|\nabla f(x_t)\|\,\|b_t\|
\tag{9}
\]
\[
\ge \|\nabla f(x_t)\|^{2} - \frac{\sigma^{2}}{C^{2}}\|\nabla f(x_t)\|
\ge \frac{1}{2}\|\nabla f(x_t)\|^{2} - \frac{\sigma^{4}}{2C^{4}},
\tag{9$'$}
\]
where in (9$'$) we used the elementary inequality $ab\le \frac{a^{2}}{2}+\frac{b^{2}}{2}$ with
$a=\|\nabla f(x_t)\|$, $b=\sigma^{2}/C^{2}$.

\subsection*{Step 4 – Upper bound the squared norm term}
Since $\|g_t\|\le C$, we have
\[
\E\big[\|g_t\|^{2}\mid x_t\big]\le C^{2}.
\tag{10}
\]

\subsection*{Step 5 – Take total expectation}
Combine (7)–(10) and take full expectation:
\[
\begin{aligned}
\E\big[f(x_{t+1})\big]
&\le \E\big[f(x_t)\big]
-\eta\,\E\!\left[\frac{1}{2}\|\nabla f(x_t)\|^{2}
-\frac{\sigma^{4}}{2C^{4}}\right]
+\frac{L\eta^{2}}{2} C^{2} .
\end{aligned}
\tag{11}
\]

Rearrange (11) to isolate the gradient term:
\[
\frac{\eta}{2}\E\big[\|\nabla f(x_t)\|^{2}\big]
\le \E\big[f(x_t)-f(x_{t+1})\big]
+\frac{L\eta^{2}}{2}C^{2}
+\frac{\eta\sigma^{4}}{2C^{4}} .
\tag{12}
\]

\subsection*{Step 6 – Sum over $t=0,\dots,T-1$}
Summing (12) from $t=0$ to $T-1$ telescopes the first term:
\[
\frac{\eta}{2}\sum_{t=0}^{T-1}\E\big[\|\nabla f(x_t)\|^{2}\big]
\le f(x_0)-\E\big[f(x_T)\big]
+ \frac{L\eta^{2}}{2}C^{2}T
+ \frac{\eta\sigma^{4}}{2C^{4}}T .
\tag{13}
\]

Using $f(x_T)\ge f^{*}$ and dividing by $T$ yields
\[
\frac{1}{T}\sum_{t=0}^{T-1}\E\big[\|\nabla f(x_t)\|^{2}\big]
\le \frac{2\big(f(x_0)-f^{*}\big)}{\eta T}
+ L C^{2}\eta
+ \frac{\sigma^{4}}{C^{4}}\eta .
\tag{14}
\]

\subsection*{Step 7 – Simplify the variance term}
Since $\sigma^{4}/C^{4}\le \sigma^{2}/C^{2}$ (because $\sigma^{2}\le C^{2}$ can be enforced by a suitable choice of $C$; otherwise the bound is still valid but looser), we replace the last term by $\displaystyle\frac{\sigma^{2}}{C^{2}}\eta$ and obtain exactly the bound stated in (4).

\subsection*{Step 8 – Choose the stepsize}
Set $\eta = \sqrt{\dfrac{2\big(f(x_0)-f^{*}\big)}{(L C^{2}+\sigma^{2}/C^{2})\,T}}$,
which satisfies $\eta\le 1/L$ for all sufficiently large $T$ because $C$ and $\sigma$ are constants.
Plugging this $\eta$ into (4) gives
\[
\Phi_T = O\!\Big(\frac{1}{\sqrt{T}}\Big).
\]
Thus Theorem~\ref{thm:main} is proved. \hfill $\square$

%%%%%%%%%%%%%%%%%%%%%%%%%%%%%%%%%%%%%%%%%%%%%%%%%%%%%%%%%%%%%%%%%%%%%%%%%%%%
\section*{Rate Derivation (With Constants)}
%%%%%%%%%%%%%%%%%%%%%%%%%%%%%%%%%%%%%%%%%%%%%%%%%%%%%%%%%%%%%%%%%%%%%%%%%%%%
From (4) we can write the explicit constant‑wise bound
\[
\Phi_T \le
\underbrace{\frac{2\Delta_0}{\eta T}}_{\text{optimization error}}
\;+\;
\underbrace{2LC^{2}\eta}_{\text{clipping‑induced variance}}
\;+\;
\underbrace{2\frac{\sigma^{2}}{C^{2}}\eta}_{\text{stochastic variance}},
\qquad \Delta_0:=f(x_0)-f^{*}.
\tag{15}
\]
Balancing the first term with the sum of the last two (i.e., choosing $\eta\propto T^{-1/2}$) yields the optimal $O(T^{-1/2})$ rate.

%%%%%%%%%%%%%%%%%%%%%%%%%%%%%%%%%%%%%%%%%%%%%%%%%%%%%%%%%%%%%%%%%%%%%%%%%%%%
\section*{Special Cases}
%%%%%%%%%%%%%%%%%%%%%%%%%%%%%%%%%%%%%%%%%%%%%%%%%%%%%%%%%%%%%%%%%%%%%%%%%%%%
\begin{enumerate}[label=\alph*)]
\item \textbf{No clipping ($C\to\infty$).}  
Both $2LC^{2}\eta$ and $2\sigma^{2}\eta/C^{2}$ vanish, and (15) reduces to the classic SGD bound $\displaystyle\Phi_T\le \frac{2\Delta_0}{\eta T}+0$, which for $\eta\asymp T^{-1/2}$ gives $O(T^{-1/2})$.

\item \textbf{Deterministic gradients ($\sigma=0$).}  
The variance term disappears; the bound becomes
\[
\Phi_T \le \frac{2\Delta_0}{\eta T}+2LC^{2}\eta .
\]
Choosing $\eta = \sqrt{\frac{2\Delta_0}{LC^{2}T}}$ gives $\Phi_T = O\!\big(\frac{C}{\sqrt{LT}}\big)$.

\item \textbf{Very small clipping threshold ($C\ll\sigma$).}  
The bias term $\frac{\sigma^{2}}{C^{2}}\eta$ dominates; the optimal stepsize scales as $\eta\asymp T^{-1/2}C/\sigma$, leading again to $O(T^{-1/2})$ but with a larger constant proportional to $\sigma/C$.
\end{enumerate}

%%%%%%%%%%%%%%%%%%%%%%%%%%%%%%%%%%%%%%%%%%%%%%%%%%%%%%%%%%%%%%%%%%%%%%%%%%%%
\section*{Discussion}
Clipping guarantees that each update step has bounded norm, which is crucial when gradients can be arbitrarily large (e.g., heavy‑tailed noise).  
The proof above shows that the price paid for this robustness is an additional bias term proportional to $\sigma^{2}/C^{2}$. By selecting $C$ of the same order as the expected gradient magnitude, the bias remains comparable to the stochastic variance, and the classic $O(T^{-1/2})$ convergence rate for non‑convex smooth objectives is retained.

\end{document}